\scanpage{28}\mspage{27}
% L1–3
\noindent\llap{\rule[-1.05\baselineskip]{0.4pt}{1.85\baselineskip}\hspace{0.65em}}\uline{Prop} Si $T,T'$ sont des ps. tr. en position rel modérée, alors $T\cap T'$ est \inserted{vide ou est} un ps. tr, et on a
\[
\partial(T\cap T')\subset\partial T\cup\partial T'
\]

% L4–6
Immédiat. On va introduire une relation d'ordre entre ps-tronçons \inserted{(pas l'inclusion!)}. On dit que
\[
T\unsure{<|\prec}T'
\]
% L7–8
si \deleted{a) $T$ est en pos. rel modérée avec $T'$,}
\[
\deleted{\text{b)}}\quad\deleted{\unsure{\forall}}\quad t\in T,\ t'\in T'\implies t<t'
\]
\margin{\leavevmode\llap{\rule[-4\baselineskip]{0.35pt}{5\baselineskip}\hspace{0.4em}}\deleted{\unsure{Si les}} Ça implique $T$ et $T'$ en pos. rel. mod et $T\cap T'=\varnothing$}
% L9
C'est \unsure{évident|évidemt} transitif.

% L10–15
Soit $T<T'$. Comme $T$ est majoré, il s'ensuit que \deleted{$\partial T\neq\varnothing$} il a un plus grand \unsure{élém|élément} \deleted{\unsure{b}} $a$, \deleted{tel que} \inserted{et alors} $T\subset L_{\leq a}$. \deleted{\unsure{Comme}} De $=$, $\partial T'$ a un plus petit élément $b$, tel que $T'\subset L_{\geq b}$, \deleted{\unsure{et on a}} \deleted{\unsure{alors}} et on voit de suite que $a\leq b$.

% L16 and note below brace
\uline{Position relative \deleted{impropre} \inserted{de raccordement}} : $a=b$, et $\underbrace{a\in T\text{ ou }b\in T'}$%
\margin{à l'exclusion l'un de l'autre, puisque $T\cap T'=\varnothing$}

% L17–20
\deleted{\unsure{Pour|Par} la suite, on l'\uline{exclut} \unsure{donc} :}

On renforce la relation $<$ en $\ll$
\[
\begin{aligned}
T\ll T' &\Longleftrightarrow T<T',\quad\text{et position relative}\\
&\hspace{12em}\text{propre}\\
&\Longleftrightarrow \exists\,x\in L\text{ tel que }T\subset L_{\unsure{<|\leq}a},\ T'\subset L_{>x}
\end{aligned}
\]
